\begin{Exercise}[label={ex:sequenza}]
  Realizzare un programma che prenda in input una sequenza di cifre (tra 1 e 9) e calcoli la somma massima fra le sotto-sequenze di cifre non decrescenti. Il programma termini quando viene inserito lo 0.
  
  Esempio:
  \begin{tabular}{|c c c c|c c|c|c c|c|}
    \hline
    2 & 2 & 4 & 5 & 3 & 9 & 3 & 1 & 5 & 0 \\
    \hline
    \textbf{13} & & & & 12 & & 3 & 6 & & \\
    \hline
  \end{tabular}

\end{Exercise}

\begin{Answer}[ref={ex:sequenza}]
  \lstinputlisting[style=cstyle]{src/sequenza.c}
\end{Answer}